% !PRINTED 277
\weberpage{277}
Uebereinstimmend damit ist auch $\sqrt{-2}$ reell oder rein imaginär
(Bd.\,I, \S.\,138, 4., 6.), und folglich ist in diesen Fällen $\Theta$,
$\psi$, $\chi$ und mithin die ganze Tetraëdergruppe reell.

2) Ist aber $p \equiv 3 \pmod 4$, so sind $\gamma^{\frac{p^2-1}{4}}$,
$\gamma^{-\frac{p^2-1}{4}}$ conjugirt imaginär, weil
$\gamma^{(p+1)\frac{p^2-1}{4}} = 1$ ist.

Ist $p \equiv 3 \pmod 8$, so ist $\sqrt{-2}$ reell und
$\gamma^{\frac{p^2-1}{8}(p+1)} = -1$, also $\gamma^{\frac{p^2-1}{8}}$
und $-\gamma^{\frac{p^2-1}{8}}$ conjugirt imaginär, und ist $p \equiv 7
\pmod 8$, so ist $\sqrt{-2}$ rein imaginär, $\gamma^{\frac{p^2-1}{8}}$
und $\gamma^{-\frac{p^2-1}{8}}$ conjugirt imaginär, und folglich gehört
$\Theta$, $\psi_1$, $\chi$ in diesen Fällen zur Gruppe $\Gamma_p$. Es
giebt also in \gesperrt{allen Fällen in der Gruppe} $L_p$ eine Tetraëdergruppe.

Für $p = 5$ z. B. kann man, da $-2$ quadratischer Nichtrest von 5 ist,
$\varepsilon = \sqrt{-2}$, und, da $\pm 2$ die beiden primitiven
Wurzeln von 5 sind, etwa
\[
\gamma^6 = -2, \quad \gamma^3 = \sqrt{-2}, \quad \gamma = 1 + \sqrt{-2}
\]
setzen, dann folgt
\weq{6}{\begin{aligned}
&\Theta = \begin{pmatrix} 2, & 0 \\ 0, & 3 \end{pmatrix},\quad
\psi = \begin{pmatrix} 0, & 2 \\ 2, & 0 \end{pmatrix},\quad
\psi_1 = \begin{pmatrix} 0, & 1 \\ -1, & 0 \end{pmatrix} \\
&\chi = \begin{pmatrix} 1, & 1 \\ 2, & -2 \end{pmatrix},\quad
\chi^2 = \begin{pmatrix} 2, & 1 \\ 2, & -1 \end{pmatrix},
\end{aligned}}
und daraus ergiebt sich die gesuchte Tetraëdergruppe für $p = 5$
[\S.\,56, (3)]:
\weq{7}{\begin{aligned}
&\begin{pmatrix} 1, & 0 \\ 0, & 1 \end{pmatrix},
\begin{pmatrix} 2, & 0 \\ 0, & -2 \end{pmatrix},
\begin{pmatrix} 0, & 2 \\ 2, & 0 \end{pmatrix},
\begin{pmatrix} 0, & 1 \\ -1, & 0 \end{pmatrix} \\
&\begin{pmatrix} 1, & 1 \\ 2, & -2 \end{pmatrix},
\begin{pmatrix} 2, & 2 \\ 1, & -1 \end{pmatrix},
\begin{pmatrix} -1, & 1 \\ 2, & 2 \end{pmatrix},
\begin{pmatrix} -2, & 2 \\ 1, & 1 \end{pmatrix} \\
&\begin{pmatrix} 2, & 1 \\ 2, & -1 \end{pmatrix},
\begin{pmatrix} -1, & 2 \\ 1, & 2 \end{pmatrix},
\begin{pmatrix} 1, & 2 \\ 1, & -2 \end{pmatrix},
\begin{pmatrix} -2, & 1 \\ 2, & 1 \end{pmatrix}.
\end{aligned}}
Diese Gruppe ist ein Theiler von $L_5$ vom Index 5.

IV. \gesperrt{Octaëdergruppe.}

In der Octaëdergruppe ist eine cyklische Gruppe vom Grade 4 enthalten,
und eine solche Gruppe kann also nur existiren, wenn $1/2\,(p-1)$ oder
$1/2\,(p+1)$ durch 4 theilbar, d. h. wenn $p \equiv \pm 1 \pmod 8$ ist.
